%%%%%%%%%%%%%%%%%%%%%%%%%%%%%%%%%
%
%       EXERCÍCIO
%
%%%%%%%%%%%%%%%%%%%%%%%%%%%%%%%%%

\ifdefstring{\atividade}{prova}{%
    \renewcommand{\valorquestao}{\ValorQ\ pontos}
}%

\begin{exercicioBanco}[\valorquestao]
O volume diário de leite produzido por uma cooperativa agrícola (em milhares de litros) foi registrado durante 40 dias consecutivos e é apresentado no rol a seguir:

\begin{center}
    \renewcommand{\arraystretch}{1.3}
    \begin{tabular}{|C{30pt}|C{30pt}|C{30pt}|C{30pt}|C{30pt}|C{30pt}|C{30pt}|C{30pt}|C{30pt}|C{30pt}|}
        \hline
        10.2 & 10.4 & 10.5 & 10.7 & 11.0 & 11.1 & 11.3 & 11.4 & 11.5 & 11.5 \\
        \hline
        11.6 & 11.6 & 11.7 & 11.7 & 11.8 & 11.8 & 11.8 & 11.9 & 11.9 & 12.0 \\
        \hline
        12.0 & 12.1 & 12.1 & 12.2 & 12.2 & 12.3 & 12.3 & 12.4 & 12.4 & 12.5 \\
        \hline
        12.6 & 12.7 & 12.8 & 13.0 & 13.1 & 13.3 & 13.5 & 13.6 & 13.8 & 14.2 \\
        \hline
    \end{tabular}
\end{center}

\begin{enumerate}[(a)]
    \item Construa o histograma dessa variável estatística, considerando intervalos de classe com amplitude igual a \(0.8\), iniciando a primeira classe em \(10.0\).
    %
    \item Construa o boxplot (diagrama de caixas) correspondente a esse conjunto de dados.
\end{enumerate}
%
\end{exercicioBanco}

